%=============================================================================!
% 17.0 CHAPTER OPENING !
%=============================================================================!
The methods of Chapters~9 to~16 are now implemented in \texttt{mdlab}.
Research codes use the same physical ideas, but their units, defaults
and data formats differ. Moving a calculation between codes therefore
starts with a comparison whose answer we already know. Here we run the
same argon model in \texttt{mdlab}, the Atomic Simulation Environment
(ASE)\cite{Larsen2017} and LAMMPS\cite{Thompson2022}, tracing the
conversions needed for their results to agree.

We then use MACE-MP-0\cite{Batatia2025}, a \emph{foundation model}
trained on first-principles calculations across much of the periodic
table. Its broad training makes it useful on many materials, but does
not establish its accuracy on a particular system. We examine its cost
and time step, and use FHI-aims\cite{Blum2009} to measure the cost and
convergence of a first-principles reference. Here \emph{first-principles}
means that the forces are calculated from the electrons
(\cref{sec:pe-origin}) rather than supplied by a fitted interatomic model.

Lithium in graphite brings these choices together in a material that the
earlier chapters represented with model surfaces and layered solids.
Its spacing and hops require different convergence checks. A production
protocol collects those checks with the measurements that justify each
setting; the final gallery links the faults of a trajectory to the
checks that reveal them.
